\chapter{Contexts, sub-agents and other models}
\label{ch:contexts}

The model sees a \emph{window} rebuilt from the log for every request. This
chapter is about shaping that window on purpose: hiding what no longer
matters, giving a helper a clean desk of its own, and handing work to a
different model.

\section{The window, briefly}

By default the window holds the system prompt (with your latest request
and its plan pinned into it, where they are never trimmed away), then the
conversation. When it grows past the \emph{high water} mark, the oldest
tool results are elided in place, down to the \emph{low water} mark. A tool
result longer than a cap is shortened. The sizes are in the config, grow
with a model whose context is larger than the defaults assume, and are
readable from Scheme:

\begin{txt}
you> :eval (window-limits)
((high-water . 61440) (low-water . 30720) (ceiling . 61440) (max-result-tokens . 3072) (output-reserve . 16384))
\end{txt}

The \code{ceiling} is the most a request may carry; above the high water it
is room kept for compaction's note. \code{(context-tokens)} is the size of
the next request, counted as the window counts.

\code{output-reserve} is also the cap on each reply. Trimming never touches
the log; it only changes what the next request contains.

\begin{tip}
The model server caches the start of the prompt between requests, and a
cached prefix is what makes the second request in a turn fast. Everything
in this chapter that changes \emph{old} parts of the window (folding an
early result, narrowing, compacting) makes the server re-read from that
point once. Do it in batches, or near the end, not on every round.
\end{tip}

\section{Folding a result}

\code{(fold-result result text)} shows a tool result as \code{text} in
every later window: an outline instead of a whole file, a line saying what
a failed command was. The result stays in the log, so you (and
\code{nodwyn replay}) still see all of it. \code{result} is the call's id,
or the \code{tool.result} event's number.

\begin{scm}
;; After a successful build, the model never needs the old build errors.
(fold-result "c7" "[folded: build errors from before the fix]")
\end{scm}

The hygiene habit (Chapter~\ref{ch:habits}) is built entirely on this.

\section{Narrowing: a clean desk}

\code{(narrow-context system [name])} starts a fresh context, like
narrowing a buffer in Emacs: until it is folded, models see \code{system}
as their system prompt and only what comes after this point. Follow it
with \code{add-message} to hand over a task.

\code{(fold-context)} ends the innermost narrowed context and folds it
away: later windows skip everything since the narrow and see the enclosing
conversation again. \code{(widen)} ends every narrowed context and shows
everything, including what happened inside them. A new user message widens
too, except for a \emph{lasting} narrowing, \code{(narrow-context system
name '((lasting . \#t)))}: if it is still open when the user speaks again,
the conversation starts there from then on. That is how a compaction
outlasts its turn.

\begin{txt}
   main conversation ..........................   (narrow-context "You review...")
                                                   | harness.msg: the task
                                                   | llm.req / llm.resp
                                                   | tool calls and results
   main conversation continues  <------------------(fold-context)
   (the reviewer's work is in the log, not in the window)
\end{txt}

Contexts nest, and calls the model made just before a narrow wait for it
to fold. You rarely need these primitives directly, because
\code{sub-agent} wraps them.

\section{Sub-agents}

\begin{txt}
(sub-agent system task tools [name])        ; [name] is optional
\end{txt}

runs \code{task} in a narrowed context under \code{system}, lets the model
use \code{tools}, and returns its final answer after folding the context
away. Whoever asked sees only the answer. \code{tools} is one of:

\begin{itemize}
  \item a list of names: the model is offered only those, and any other
    call is skipped;
  \item a procedure that takes a call and says whether it may run (the
    model is then offered every tool);
  \item \code{\#f} for all tools.
\end{itemize}

\subsection{A reviewer at the end of every turn}

\begin{scm}
(define review-base-agent-turn agent-turn)
(define (agent-turn request)
  (review-base-agent-turn request)
  (let ((verdict (sub-agent "You review a summary of a folder. Reply OK or name what is missing."
                            (string-append "The request was: " request)
                            '("fs-list" "fs-read")
                            "reviewer")))
    (log-note (string-append "reviewer: " verdict))))
\end{scm}

And its trail in the log:

\begin{txt}
   13 context.narrow  {System:You review a summary ... Agent:reviewer Answers:[]}
   14 harness.msg     {Text:The request was: Summarize this folder}
   15 llm.req         {... Messages:2 Tools:2 ...}
   16 llm.resp        {Content:OK ... DoneReason:stop ...}
   17 context.fold    {}
   18 note            {Text:reviewer: OK}
\end{txt}

\code{Messages:2 Tools:2}: the reviewer saw its own system prompt and task,
and was offered exactly two tools. None of the main conversation.

\subsection{Agent tools: a sub-agent the model can call}

\begin{scm}
(define-agent-tool "ask-gardener"
  "Ask an experienced gardener, who can read the garden's files, a question."
  '(("question" "string" "what you want to know"))
  (lambda (call args)
    (sub-agent "You are an experienced gardener. Read what you need, then answer in two sentences."
               (assoc-ref args "question")
               '("fs-read" "fs-list"))))
\end{scm}

The model sees \code{ask-gardener} among its tools. When it calls it,
\code{run-calls} runs your handler instead of a tool body and answers the
call with \code{answer-call}. Define agent tools in \code{init.scm}, before
the tool list is fixed.

\subsection{Sub-agents stop themselves, not you}

A sub-agent that loops would trip the core's progress stop and end your
whole turn. So \code{sub-agent} stops short of it on its own: an identical
call more than \code{sub-agent-repeat-limit} (2) times in a row is skipped,
and if it repeats again the sub-agent is asked for its answer without
tools, which comes back with a note that it may be incomplete. After
\code{sub-agent-max-rounds} (40) rounds it is likewise asked to wrap up. Its
caller's work goes on either way.

\section{Compaction}

A long task fills the window with tool output the model no longer needs.
Past \code{compact-threshold} tokens of \code{(context-tokens)} (by default
just under the high water: 58,983 of 61,440), \code{maybe-compact} asks the model for a
\textbf{handover note}: the task and its example as given, what it found
with the evidence for each finding, what it changed, how to test, what is
unsure, and its check script. Then it narrows the context to the system
prompt, the request and that note, and the work continues, from a clean
window, with the changes already in the files. The narrowing lasts: the
next turns start from the note too, not from the whole conversation.

You can compact whenever you like, whatever the size, by sending a message:

\begin{txt}
/compact                     compact now
/compact the failing tests   and keep these in particular
\end{txt}

The note is written for your latest request before the \code{/compact},
and the turn ends there: your next message goes on from the note. A
strategy asks for the same at its next round with
\code{(compact-next-round)}.

Every compaction is one line wherever you watch, and one
\code{context.compact} record in the log: what set it off, and how big the
context was before and after, in the window's estimate (the unit of the
threshold; your provider usually counts a fifth or so fewer):

\begin{txt}
── compacted (past ~58,983 tokens): context ~59,210 → ~4,100 tokens, handover note 812 tokens ──
\end{txt}

One that gave up, because no note came back, is a warning, and the turn
goes on in the full context. How full the context is shows as you work:
in the terminal client's title (\code{context \textasciitilde 45k of
\textasciitilde 61k}) and \code{:budget}, and in the console's Context
meter.

\begin{scm}
(define compact-threshold 40000)   ; compact earlier
(define compact-threshold #f)      ; never compact; only trim
\end{scm}

A strategy that replaces \code{agent-turn} should call
\code{(maybe-compact request reply)} where the default does: after a
reply's calls are answered, before asking again.

\begin{careful}
Compaction is a trade. Measured on a coding benchmark, compacting at
\emph{half} the window stopped runs that would have finished, to write
notes and re-orient from them (\adr{32}). The late default compacts only
work that really runs long.
\end{careful}

\section{Other models}

Any model you configure can be asked by name:

\begin{toml}
# config.toml
[models.coder]
provider = "openai"
url = "https://api.openai.com/v1"
model = "gpt-5.1-codex"
credential = "openai-key"          # a keychain class, bound to api.openai.com

[models.local]
provider = "ollama"
url = "http://localhost:11434"
model = "qwen3:8b"
\end{toml}

\begin{scm}
(model-names)                ;=> ("coder" "control" "local")
(model-info "local")         ;=> ((name . "local") (model . "qwen3:8b") (provider . "ollama")
                             ;    (num-ctx . 77824) (max-output . 16384)
                             ;    (temperature . #f) (seed . #f))
(model-turn "coder")         ; inside a turn: ask the coder
\end{scm}

\subsection{Sampling: repeatable replies}

Any model can carry its own sampling settings. Unset, the server decides:

\begin{toml}
[models.local]
provider = "ollama"
url = "http://localhost:11434"
model = "qwen3:8b"
temperature = 0          # 0 to 2
seed = 42
\end{toml}

They go with every request to that model and are recorded with it, so
\code{nodwyn replay -{}-window} still rebuilds the request exactly. A
temperature of 0 with a seed makes replies as repeatable as the server
allows. That is useful when you compare two strategies on the same task,
though many servers are not bit-for-bit deterministic across hardware or
batching. \code{model-info} shows them, as \code{\#f} when unset.

\textbf{Routing} is narrowing plus naming: give a model a context you
built, and let it work.

\begin{scm}
;; The control model plans; the coder implements from the plan alone.
(define (implement-with-coder plan)
  (narrow-context (session-system-prompt) "coder")
  (add-message (string-append "Implement this plan, then report what you did:\n" plan))
  (let loop ((reply (model-turn "coder")))
    (let ((calls (assoc-ref reply 'tool-calls '())))
      (if (null? calls)
          (begin (fold-context) (assoc-ref reply 'content ""))
          (begin (run-calls calls)
                 (loop (model-turn "coder")))))))
\end{scm}

\begin{careful}
Configuring a model is consenting that what a strategy sends it leaves
your machine. With \code{narrow-context} that can be just a task and a
plan; a strategy that asks a remote model over the whole conversation sends
the whole conversation. The key goes only where its binding says, and only
you configure models: the model cannot.
\end{careful}

\begin{puzzle}[ 16: second opinion]
Write an agent tool \code{second-opinion} that asks the model named
\code{"local"} the same question with no tools at all, in a narrowed
context, and returns its answer. (Hint: \code{model-turn} takes options.)\ans{16}
\end{puzzle}
